\documentclass[12pt]{article}

\usepackage{geometry}
\geometry{a4paper, portrait, margin=1in}

\usepackage{setspace}
\doublespacing


%need to format subitems

\begin{document}
\title{Brazil Chapter}
\author{Kenyon P. Cavender}
\maketitle

\section{Introduction}
Until 1913, the Amazon was by far the largest producer of rubber for the world market. During the Amazon Rubber Boom a massive influx of foreign capital transformed life in the Amazon. The development of infrastructure connected the relatively isolated region to global markets, and rubber barons subjected the people of the Amazon, as well as displaced migrants from other regions of Brazil looking for work, to new forms of labor exploitation. 

\section{The Free Laborer}
This section looks at one sense of the freedom of proletarianization - the dissolution of traditional social ties and community relationships. In the nineteenth century, three major events undermined the relations of dependency that structured Brazilian society. They were the \textit{Cabanagem}, the \textit{Grand Seca}, and the abolition of slavery. Each of these events set loose some portion of the population which had previously been bound in a customary social hierarchy. Unlike their fellow workers of the North, however, they did not become proletarians, selling their labor for a wage and meeting their needs through a market of commodities. They were rather cast adrift to scratch out a living as \textit{seringueiros}, or rubber workers, in the Amazon where market penetration was both shallow and mediated by the workers' debt contracts.
\end{document}
